% =============================================================
% integration/chapter1.tex
% Integration wrapper for Chapter 1 (FROZEN SOURCE).
%
% This file adapts the standalone Chapter1/main.tex for inclusion
% in the dissertation master document without modifying the frozen source.
%
% Frozen source: Chapter1/main.tex
% Standalone commands suppressed: \documentclass, \begin{document},
%   \maketitle, \end{document}
% Chapter-opening material reproduced from frozen source content.
%
% INTEGRATION LEDGER: see docs/03_INTEGRATION_LEDGER.md
% =============================================================

% ---- Set graphics search path for this chapter --------------
% \subimport changes the current dir, so figures referenced as
% "figures/foo.pdf" will resolve to Chapter1/figures/foo.pdf
\graphicspath{%
  {Chapter1/}%
  {Chapter1/figures/}%
  {Chapter1/appendixA/figures/}%
}

% ---- Chapter opening (replaces standalone \maketitle) -------
% Content extracted verbatim from Chapter1/main.tex title block.
% Suppressed: \thanks{} footnotes for author affiliation and AI note
%   (dissertation title page identifies the author globally).
%   The AI disclosure note is preserved in the footnote below.

\chapter[Essay 1 \textemdash{} Critical Replication of Shaikh's Capacity Utilization Measure]%
  {Critical Replication of Shaikh's Capacity Utilization Measure}

\vspace{0.8em}

% Abstract (verbatim from frozen source)
\begin{abstract}
Identifying a productive capacity ceiling requires estimating an unobserved
benchmark, a task that risks conflating technical engineering limits with the
institutional realities of capitalist accumulation. This chapter critically
replicates Anwar Shaikh's accounting-based capacity measure for the post-war
US corporate sector (1947--2011), reinterpreting the long-run output--capital
coefficient as a capacity transformation elasticity ($\theta$) under conditions
of unbalanced growth. Testing this relation across a combinatorial grid of 500
single-equation specifications and a system-level Vector Error Correction Model
(VECM), we demonstrate that the standalone bivariate output-capital relation
fails to cointegrate. System-level stability emerges exclusively when the rate
of exploitation ($e_t$) enters the cointegrating vector, yielding an estimated
transformation elasticity ($\hat{\theta} < 1$) indicative of structural
overaccumulation. This evidence suggests that the productive ceiling is not a
neutral engineering constraint. Rather, it is a contingent outcome of
decentralized accumulation, structurally conditioned by the distributive
conflict and the labor process within and beyond the workplace, which physically
dictate the conversion of accumulated capital into active productive capacity.
\end{abstract}

\vspace{0.8em}

\noindent\textbf{Keywords:} Capacity Utilization; Critical Replication; Shaikh;
Transformation Elasticity; Cointegration; ARDL; VECM; Distribution.

\noindent\textbf{JEL Codes:} B51; C32; E01; E22; E32; P16.

\clearpage

% ---- Chapter body (frozen source sections) ------------------
% \subimport{Chapter1/}{file} sets current dir to Chapter1/ so that
% any nested \input{} calls within section files resolve correctly.
\subimport{Chapter1/}{section1}
\subimport{Chapter1/}{section2}
\subimport{Chapter1/}{section3}
\subimport{Chapter1/}{section4}
\subimport{Chapter1/}{section5}

% ---- Chapter bibliography -----------------------------------
\clearpage
{\small
\bibliographystyle{apalike}
\bibliography{Chapter1/references}
}

% ---- Chapter appendices -------------------------------------
\clearpage
\begin{subappendices}
% AppendixODE has no sub-inputs; either base dir works.
\subimport{Chapter1/}{AppendixODE/AppendixODE}
\clearpage
% appendixA.tex uses \input{appendixA/tables/...}, so current dir
% must be Chapter1/ (not Chapter1/appendixA/) for paths to resolve.
\subimport{Chapter1/}{appendixA/appendixA}
\end{subappendices}
